% Báo cáo tiến độ Đồ án 1 (Gồm trang bìa chuẩn HUST và nội dung gói gọn đúng trong 1 trang A4)
% Định dạng tuân thủ chặt chẽ quy chuẩn trình bày báo cáo ĐHBK Hà Nội (report/)
\documentclass[11pt,a4paper]{article}

\usepackage{fontspec}
\usepackage[vietnamese]{babel}
\setmainfont{Times New Roman}
\setmonofont{Consolas}[Scale=0.85]

\usepackage[
  a4paper,
  top=1.85cm,
  bottom=1.85cm,
  left=2.5cm,
  right=2.0cm,
  headheight=14pt
]{geometry}

\usepackage{setspace}
\setstretch{1.22}
\setlength{\parskip}{3.5pt}

\usepackage{graphicx}
\usepackage{booktabs}
\usepackage{tabularx}
\usepackage{array}
\usepackage{enumitem}
\usepackage{tikz}
\usetikzlibrary{calc}
\usepackage{titlesec}
\usepackage{indentfirst}
\setlength{\parindent}{1cm}
\usepackage{fancyhdr}
\usepackage[hidelinks]{hyperref}

% Định dạng tiêu đề các mục theo chuẩn học thuật của report/
\titleformat{\section}
  {\normalfont\bfseries\fontsize{13}{16}\selectfont}
  {\thesection.}{0.45em}{}
\titlespacing*{\section}{0pt}{12pt}{5pt}

% Danh sách căn chỉnh trang nhã, khoảng cách vừa vặn
\setlist[itemize]{itemsep=4pt,topsep=3pt,leftmargin=1.5em,parsep=0pt}

% Thiết lập số trang đáy trang giữa theo chuẩn báo cáo
\pagestyle{fancy}
\fancyhf{}
\fancyfoot[C]{\fontsize{10}{12}\selectfont\thepage}
\renewcommand{\headrulewidth}{0pt}
\renewcommand{\footrulewidth}{0pt}

\begin{document}

% ================================================================
% TRANG 1: TRANG BÌA (Chuẩn mẫu bìa ĐHBK Hà Nội / Trường Điện - Điện tử)
% ================================================================
\newgeometry{top=2cm,bottom=2cm,left=2cm,right=2cm}
\begin{titlepage}
\thispagestyle{empty}
\begin{tikzpicture}[remember picture,overlay]
  \draw[line width=1.0pt]
    ($(current page.north west)+(1.55cm,-1.35cm)$)
    rectangle
    ($(current page.south east)+(-1.55cm,1.35cm)$);
  \draw[line width=0.4pt]
    ($(current page.north west)+(1.63cm,-1.43cm)$)
    rectangle
    ($(current page.south east)+(-1.63cm,1.43cm)$);
\end{tikzpicture}
\begin{center}
  {\bfseries\fontsize{20}{25}\selectfont
  ĐẠI HỌC BÁCH KHOA HÀ NỘI\\
  TRƯỜNG ĐIỆN -- ĐIỆN TỬ\par}

  \vspace{0.2cm}
  \rule{6cm}{0.4pt}

  \vspace{0.9cm}
  \includegraphics[height=4.6cm]{figure/logo.png}

  \vspace{0.8cm}
  {\bfseries\fontsize{20}{25}\selectfont BÁO CÁO TIẾN ĐỘ\\[0.25cm]
  ĐỒ ÁN 1\par}

  \vspace{1.0cm}
  {\bfseries\fontsize{16}{22}\selectfont
  HỆ THỐNG CHATBOT TƯ VẤN CHỌN MÔN HỌC\\CHO SINH VIÊN\par}

  \vfill

  {\fontsize{13}{18}\selectfont
  \begin{tabular}{@{}>{\bfseries}l@{\hspace{1.2em}}l@{}}
    Giảng viên hướng dẫn: & TS. Bùi Tiến Cường \\
    Sinh viên thực hiện:  & Mai Duy Hiếu \\
  \end{tabular}\par}

  \vfill
  {\bfseries\fontsize{13}{17}\selectfont Hà Nội, 10/2026}
\end{center}
\end{titlepage}
\restoregeometry
\clearpage

% ================================================================
% TRANG 2: NỘI DUNG TIẾN ĐỘ (Đúng chuẩn layout report/, gói gọn 1 trang)
% ================================================================
\setcounter{page}{1}

\begin{center}
  {\bfseries\fontsize{15}{18}\selectfont BÁO CÁO TIẾN ĐỘ ĐỒ ÁN 1}\par
\end{center}
\vspace{0.35cm}

\section{Tiến độ}

\begin{itemize}
  \item \textbf{Cải thiện hệ thống RAG và chất lượng phản hồi của Chatbot:} Chuẩn hóa kho dữ liệu vector bằng các văn bản kiểm chứng chính thức gồm Sổ tay sinh viên 2026, Quyết định Học phí 2025--2026, chương trình đào tạo ngành Kỹ thuật Điện tử Viễn thông khóa 67 và khóa 68; hoàn thiện bộ quy tắc kiểm tra điều kiện học phần (tiên quyết, học trước) trước khi đưa ra gợi ý môn học; tích hợp cơ chế trích xuất nguyên văn có dẫn nguồn tham khảo khi dịch vụ mô hình ngôn ngữ gặp sự cố.
  
  \item \textbf{Xây dựng bộ dữ liệu đánh giá định lượng:} Thiết lập bộ dữ liệu chuẩn gồm 96 câu hỏi kiểm thử thuộc 9 nhóm tác vụ thực tế (tra cứu quy chế, học phí, học bổng, kiểm tra môn học, rà soát tốt nghiệp, gợi ý môn học, mô phỏng điểm tích lũy...); hoàn thành quy trình kiểm thử tự động đo lường định lượng các chỉ số truy xuất tài liệu, độ chính xác phân loại ý định, độ chính xác nhận diện mã môn học và thời gian phản hồi thay vì đánh giá cảm tính.
  
  \item \textbf{Chuyển đổi cơ sở dữ liệu sang MongoDB:} Chuyển đổi toàn diện hệ thống lưu trữ sang cơ sở dữ liệu hướng tài liệu MongoDB nhằm quản lý linh hoạt hồ sơ sinh viên, thông tin cá nhân và bảng điểm lồng nhau; lưu trữ toàn bộ lịch sử hội thoại của người dùng; cấu hình chỉ mục tự động thu hồi để dọn dẹp các phiên làm việc hết hạn; xây dựng kịch bản chuyển đổi dữ liệu tự động và hỗ trợ cơ chế lưu trữ dự phòng khi chạy ngoại tuyến.
  
  \item \textbf{Cải thiện giao diện và bổ sung tính năng người dùng:} Thiết kế bộ nhận diện thương hiệu gồm biểu trưng ứng dụng và bảng màu chuẩn Đại học Bách Khoa Hà Nội; bổ sung màn hình tổng hợp nhanh điểm tích lũy, số tín chỉ và cảnh báo học tập; tích hợp chức năng nhập bảng điểm trực tiếp để hệ thống tự động rà soát tiến độ tốt nghiệp và đưa ra tư vấn môn học tức thời; tối ưu hóa thanh điều hướng và menu quản lý tài khoản.
\end{itemize}

\section{Kế hoạch tiếp theo}

\begin{itemize}
  \item \textbf{Mở rộng nguồn tài liệu tham khảo:} Tiếp tục thu thập và bổ sung nhiều nguồn tài liệu tham khảo chính thức từ nhà trường, bao gồm các văn bản quy định đào tạo mới, quy chế công tác sinh viên, sổ tay sinh viên các khóa mới và khung chương trình đào tạo của các ngành đào tạo khác trong toàn trường.
  
  \item \textbf{Phát triển tính năng Đăng ký, Đăng nhập và Bảo mật:} Nghiên cứu và hiện thực hóa phân hệ quản lý tài khoản người dùng hoàn chỉnh gồm chức năng Đăng ký tài khoản, Đăng nhập và Đổi mật khẩu.
  
  \item \textbf{Hoàn thiện giao diện người dùng và tối ưu hoạt ảnh:} Hoàn thiện toàn diện thiết kế giao diện UI/UX trên các màn hình chức năng; tối ưu hóa các hiệu ứng hoạt ảnh chuyển động mượt mà cho chatbot trò chuyện.
  
  \item \textbf{Tối ưu hóa chất lượng tư vấn:} Nâng cao độ chính xác của khâu phân đoạn tài liệu quy chế đối với các điều khoản có bảng biểu phức tạp; tinh chỉnh bộ lọc và xếp hạng lại văn bản truy xuất.
\end{itemize}

\end{document}
